% Solutions to the exercises of Chapter 1.
\begin{solutions}
  \item With $b_a = 5~\mathrm{mg} = 4.905\times10^{-2}~\mathrm{m/s^2}$, the position error $\tfrac12b_at^2$ reaches 1~m at
        \[
          t = \sqrt{2/b_a} = 6.39~\mathrm{s}.
        \]
        A gyro bias produces $\tfrac16gb_gt^3$, which equals 1~m at the same time when
        $b_g = 6/(gt^3) = 3b_a/(gt) = 2.35\times10^{-3}~\mathrm{rad/s}$, about $0.13~^\circ/\mathrm{s}$ or $480~^\circ/\mathrm{h}$.
        The gyro term grows like $t^3$ and the accelerometer term like $t^2$, so for longer flights the gyro dominates: a
        typical MEMS gyro bias of $10^{-3}$~rad/s overtakes the 5~mg accelerometer bias after $t = 3b_a/(gb_g) = 15$~s.

  \item Let the vehicle be level and at rest, and let the estimate be tilted by $\delta\theta_0$ about the $x$-axis,
        $\hat R = R_x(\delta\theta_0)$. The accelerometer measures the specific force $f = (0, 0, g)$, and the navigation
        equations compute the acceleration
        \[
          \hat Rf + g_w = \big(0,\ -g\sin\delta\theta_0,\ g\cos\delta\theta_0\big) + (0, 0, -g)
          = \big(0,\ -g\sin\delta\theta_0,\ -g(1 - \cos\delta\theta_0)\big).
        \]
        The horizontal part, of magnitude $g\sin\delta\theta_0 = g\,\delta\theta_0 + O(\delta\theta_0^3)$, is the gravity leak; the
        vertical part $-g(1 - \cos\delta\theta_0) \approx -\tfrac12g\,\delta\theta_0^2$ is second order. (The exact magnitude of the
        whole error vector is $2g\sin(\delta\theta_0/2)$.) The first-order formula overestimates the horizontal leak by the
        factor $\delta\theta_0/\sin\delta\theta_0 \approx 1 + \delta\theta_0^2/6$, which reaches $1.01$ at
        $\delta\theta_0 = 0.244~\mathrm{rad} \approx 14^\circ$. For the milliradian tilts of a working filter, $g\,\delta\theta$ is
        essentially exact; it becomes inaccurate only for gross initialization errors.

  \item A camera $i$ with pose $(R_i, p_i)$ sees the point $p_f$ at camera coordinates $R_i\T(p_f - p_i)$. After the common
        transformation $R_i \mapsto RR_i$, $p_i \mapsto Rp_i + t$, $p_f \mapsto Rp_f + t$,
        \[
          (RR_i)\T\big(Rp_f + t - Rp_i - t\big) = R_i\T R\T R\,(p_f - p_i) = R_i\T(p_f - p_i),
        \]
        so every image is unchanged. Combined with \cref{prop:scale-ambiguity}, the images are invariant under every similarity
        transformation $p \mapsto sRp + t$: a monocular camera has $3 + 3 + 1 = 7$ unobservable degrees of freedom (rotation,
        translation and scale). An IMU removes three of them. The accelerometer measures acceleration in metric units, and
        scaling the trajectory by $s$ scales its accelerations by $s$, so scale becomes observable whenever the body
        accelerates. The accelerometer also measures the direction of gravity in the body frame, which a rotation about a
        horizontal axis would change, so roll and pitch become observable. Rotations about the vertical and translations
        change neither gravity nor any acceleration, so four unobservable directions remain: global position and yaw
        (\cref{thm:vio-unobservable}).

  \item During a perfect hover the true rate, velocity and acceleration are all zero. Consider what each sensor reports.
        \begin{itemize}[leftmargin=1.2em]
          \item The \emph{camera} sees identical images. The relative rotation between frames is observed to be zero, but
                with no baseline there is no parallax: neither the depths of the features nor any translation can be
                measured, because an unchanged image is also what a slowly translating camera sees in a distant scene.
          \item The \emph{gyroscope} reads $b_g$ plus noise. Since the camera pins the true rotation to zero, the reading is
                the bias: $b_g$ is estimable, exactly as in the static calibration of \cref{ch:strapdown}.
          \item The \emph{accelerometer} reads $R_{wb}\T(-g_w) + b_a$. Its direction gives roll and pitch, but only up to the
                horizontal bias: a horizontal $b_a$ is indistinguishable from a tilt of $\|b_a\|/g$, because without
                rotation the bias and gravity never move relative to each other.
        \end{itemize}
        Hence the gyro bias, and roll and pitch combined with the horizontal accelerometer bias, can be estimated. What
        cannot be distinguished: tilt from horizontal $b_a$; a velocity error, or a constant acceleration error such as the
        vertical bias, from the true (zero) motion, since nothing measures translation and such errors integrate unseen;
        the metric scale and the feature depths (no acceleration, no baseline); and, as always, global position and yaw.
        This is the ``at rest'' row of the degenerate motions of \cref{ch:observability}. A real system detects the hover and
        applies a zero-velocity update (\cref{ch:updates}): the pseudo-measurement $v = 0$ stops the velocity and position
        drift and makes the vertical bias observable, while tilt and horizontal bias separate only once the vehicle
        rotates.

  \item A model answer: \emph{EKF-SLAM puts every observed feature into the state vector, so the cost of each update grows
        quadratically with the number of features, which makes it impossible to use the hundreds of features visible in
        every image in real time. The MSCKF instead keeps a sliding window of past camera poses in the state and turns each
        feature track into constraints among all the poses that observed it, eliminating the feature position from the
        linearized equations, so the cost becomes linear in the number of features while the geometric information of
        each track is still used.}
\end{solutions}
